\section{Methodology}
\label{sec:methodology}

\subsection{Algorithm}
\label{sec:methodology:algorithm}

We apply the minimum-edge-cut recursive bisection algorithm of
\citet{dellaLibera2026recursive} via the \texttt{redist} Rust binary.
The algorithm partitions a state's census-unit adjacency graph
$G = (V, E, w)$ into $k$ balanced parts minimising total edge-cut weight,
where edge weights $w(e)$ are the TIGER boundary lengths between adjacent units.
Recursive bisection proceeds by repeatedly halving the current partition
until $k$ districts are obtained; for non-power-of-2 values of $k$, the
algorithm uses the nearest-power-of-2 strategy of \citet{dellaLibera2026nwayRecursive}.

The implementation uses METIS~5.x \citep{karypis1998fast} with
\texttt{ufactor}=5 (0.5\% per-part imbalance target) and a fixed
random seed for reproducibility. All 50-state results use seed~42
as the primary seed, with 5-seed sensitivity checks for selected states.

\subsection{Resolution Selection}
\label{sec:methodology:resolution}

We apply the resolution selection rule derived in companion paper F.3
\citep{dellaLibera2026stateLegResolution}:
\begin{itemize}
  \item \textbf{Block-group resolution} (\texttt{--resolution block\_group})
    when $k / n_{\mathrm{tracts}} > 0.05$, where $n_{\mathrm{tracts}}$ is
    the 2020 census tract count for the state.
  \item \textbf{Tract resolution} (\texttt{--resolution tract}) otherwise.
\end{itemize}

Table~\ref{tab:resolution-assignment} shows the resolution assignment for
all 50 states. Of the 50 lower chambers, 12 require block-group resolution.
Five states (NH, WY, VT, SD, and ND) have $k/n > 0.5$ and in principle
require block-level redistricting; for this paper we use block-group
resolution for these states and flag the limitation.

\begin{table}[h]
\centering\small
\begin{tabular}{llrrrr}
\toprule
State & Chamber & $k$ & $n_{\rm tracts}$ & $k/n$ & Resolution \\
\midrule
New Hampshire & House & 400 & 321 & 1.25 & block\_group \\
Wyoming       & House & 60  & 139 & 0.43 & block\_group \\
Vermont       & House & 150 & 186 & 0.81 & block\_group \\
South Dakota  & House & 70  & 183 & 0.38 & block\_group \\
North Dakota  & House & 94  & 205 & 0.46 & block\_group \\
Alaska        & House & 40  & 175 & 0.23 & block\_group \\
Montana       & House & 100 & 278 & 0.36 & block\_group \\
Nebraska      & Legislature & 49 & 570 & 0.086 & block\_group \\
Washington    & House & 98  & 1{,}458 & 0.067 & block\_group \\
Idaho         & House & 70  & 298 & 0.23 & block\_group \\
Maine         & House & 151 & 358 & 0.42 & block\_group \\
West Virginia & House & 100 & 484 & 0.21 & block\_group \\
\midrule
California    & Assembly & 80  & 8{,}997 & 0.009 & tract \\
Texas         & House & 150 & 5{,}265 & 0.028 & tract \\
Florida       & House & 120 & 4{,}245 & 0.028 & tract \\
New York      & Assembly & 150 & 4{,}919 & 0.030 & tract \\
(36 other states) & House & --- & --- & $<0.05$ & tract \\
\bottomrule
\end{tabular}
\caption{Resolution assignment for selected states. States with
  $k/n_{\rm tracts} > 0.05$ use block-group resolution (12 states total);
  the remainder use tract resolution.}
\label{tab:resolution-assignment}
\end{table}

\paragraph{Data prerequisites.}
Block-group adjacency data for the 12 block-group-resolution states must
be built before running the pipeline:
\begin{verbatim}
python scripts/data/geography/build_unit_adjacency.py \
  --resolution block_group --year 2020 \
  --states NH WY VT SD ND AK MT NE WA ID ME WV
\end{verbatim}
This generates approximately 800~MB of \texttt{.adj.bin} adjacency files.
The script is the standard adjacency builder from the F.3 data pipeline;
F.1 depends on it as a prerequisite.

\subsection{Compactness Metrics}
\label{sec:methodology:metrics}

We report Polsby-Popper compactness \citep{polsby1991third}:
\[
  PP_i = \frac{4\pi A_i}{P_i^2}
\]
where $A_i$ is the area of district $i$ and $P_i$ is its perimeter.
Values range from 0 (maximally non-compact) to 1 (perfectly circular).
We report mean PP ($\overline{PP}$), minimum PP ($\min PP$), and the
fraction of districts with $PP < 0.20$ (the ``non-compact'' threshold
used in several state legal standards).

\subsection{Population Balance}
\label{sec:methodology:balance}

Ideal district population for a state with total population $T$ and $k$ seats
is $T^* = T/k$. Population deviation for district $i$ is
$\delta_i = (T_i - T^*)/T^*$. We report maximum absolute deviation
$\max_i |\delta_i|$ and the fraction of districts within $\pm 0.5\%$.

\subsection{County Splits}
\label{sec:methodology:splits}

A county split occurs when a county's census units are assigned to more
than one district. We count the number of county splits for each state's
algorithmic map and compare to the enacted 2020 state house map where available.
